\documentclass[../../../include/open-logic-section]{subfiles}

\begin{document}

\olfileid[fr]{sth}{ord-arithmetic}{add}
\olsection{Addition ordinale}

Supposons que nous voulions additionner $\alpha$ et $\beta$.
Il suffit de placer une {copie} de $\beta$ immédiatement après
une copie de $\alpha$. (Il nous faut des \emph{copies}, puisque
nous savons, par la \olref[ordinals][basic]{ordinalsaresubsets},
que soit $\alpha \subseteq \beta$, soit
$\beta \subseteq \alpha$.) Intuitivement, nous parcourons
une suite de $\alpha$ étapes, \emph{puis} une suite de
$\beta$ étapes. La suite obtenue est bien ordonnée ;
d'après le
\olref[ordinals][ordtype]{thmOrdinalRepresentation},
elle est donc isomorphe à un ordinal (unique). Cet ordinal
peut être considéré comme la \emph{somme} de $\alpha$
et de $\beta$.

Telle est l'idée intuitive de l'addition ordinale. Pour la
définir rigoureusement, commençons par la notion de
\emph{copies} d'ensembles. Nous employons les étiquettes
arbitraires $0$ et $1$ pour indiquer la provenance
de chaque objet :

\begin{defn}\ollabel{defdissum}
La \emph{somme disjointe} de $A$ et $B$ est
$A \disjointsum B = (A\times \{0\}) \cup
(B \times \{1\})$.
\end{defn}

Définissons maintenant un ordre sur les paires d'ordinaux :

\begin{defn}
Pour tous ordinaux $\alpha_1, \alpha_2, \beta_1, \beta_2$,
posons :
\begin{align*}
	\tuple{\alpha_1, \alpha_2} \rlexless
	\tuple{\beta_1, \beta_2}\text{ ssi }&
	\text{soit $\alpha_2 \in \beta_2$}\\
	& \text{soit $\alpha_2 = \beta_2$ et $\alpha_1 \in \beta_1$}
\end{align*}
\end{defn}

Il s'agit d'un ordre \emph{lexicographique inverse} :
on compare d'abord la seconde composante, puis la première.
Rappelons que nous voulions définir
$\alpha \ordplus \beta$ comme le type d'ordre d'une copie
de $\alpha$ suivie d'une copie de $\beta$. Posons donc :

\begin{defn}\ollabel{defordplus}
Pour tous ordinaux $\alpha$ et $\beta$, leur somme est
$\alpha \ordplus \beta =
\ordtype{\alpha \disjointsum \beta, \rlexless}$.
\end{defn}
\noindent
Notons que nous avons ici légèrement abusé de la notation :
en toute rigueur, il faudrait écrire
«~$\Setabs{\tuple{x,y}\in (\alpha\disjointsum\beta)\times(\alpha\disjointsum\beta)}
{x \rlexless y}$~» à la place de «~$\rlexless$~».\footnote{Note de l'édition française : une relation sur la somme disjointe est un ensemble de paires de ses éléments ; le produit cartésien manque dans la borne source et est rétabli.}
Pour alléger l'écriture, nous continuerons d'employer
cette notation dans la suite.

Le résultat suivant, joint au
\olref[ordinals][ordtype]{thmOrdinalRepresentation},
confirme que notre définition est correcte :

\begin{lem}\ollabel{ordsumlessiswo}
Pour tous ordinaux $\alpha$ et $\beta$,
$\tuple{\alpha \disjointsum \beta, \rlexless}$
est un bon ordre.
\end{lem}

\begin{proof}
La relation $\rlexless$ est manifestement connexe sur
$\alpha \disjointsum \beta$. Pour montrer qu'elle est
bien fondée, fixons un ensemble non vide
$X \subseteq \alpha \disjointsum \beta$.
Soit $Y$ le sous-ensemble de $X$ dont les éléments
ont la plus petite seconde coordonnée possible,
c'est-à-dire
$Y = \Setabs{\tuple{\gamma, i} \in X}
{(\forall \tuple{\delta, j} \in X)i \leq j}$.
Choisissons ensuite dans $Y$ un élément dont
la première coordonnée est minimale.
\end{proof}
\noindent
Nous avons donc une définition explicite de l'addition
ordinale. Voici un fait peu surprenant (rappelons que
$1 = \{0\}$, d'après la \olref[z][infinity-again]{defnomega}) :
\begin{prop}
Pour tout ordinal $\alpha$,
$\alpha \ordplus  1 = \ordsucc{\alpha}$.
\end{prop}

\begin{proof}
Considérons l'isomorphisme $f$ de
$\ordsucc{\alpha} = \alpha \cup \{\alpha\}$ vers
$\alpha\disjointsum1 = (\alpha \times \{0\})
\cup (\{0\} \times \{1\})$, défini par
$f(\gamma) = \tuple{\gamma, 0}$ pour
$\gamma \in \alpha$, et
$f(\alpha) = \tuple{0, 1}$.\footnote{Note de l'édition
française : l'original met une seconde somme disjointe
entre les deux ensembles de paires dans cette égalité ;
la définition de $\alpha\disjointsum1$ exige leur union.}
\end{proof}
\noindent
Il est également facile de montrer que l'addition
satisfait certaines équations de récursion :

\begin{lem}\ollabel{ordadditionrecursion}
Pour tous ordinaux $\alpha, \beta$, on a :
\begin{align*}
	\alpha\ordplus 0 &= \alpha\\
	\alpha \ordplus  (\beta\ordplus 1)
	&= (\alpha \ordplus  \beta) \ordplus  1\\
	\alpha  \ordplus  \beta
	&= \supstrict_{\delta < \beta}(\alpha \ordplus  \delta)
	&& \text{si $\beta$ est un ordinal limite}
\end{align*}
\end{lem}

\begin{proof}
Vérifions les différents cas. D'abord :
\begin{align*}
	\alpha \ordplus  0
	& = \ordtype{(\alpha \times \{0\})
	\cup (0 \times \{1\}), \rlexless} \\
	&= \ordtype{(\alpha \times \{0\})
	\cup \emptyset, \rlexless}\\
	&= \alpha\\
	\alpha \ordplus (\beta \ordplus  1)
	%&= \ordtype{(\alpha\times \{0\}) \cup (\ordtype{\beta \ordplus  1}\times \{1\}), \rlexless} \\
	&= \ordtype{(\alpha\times \{0\})
	\cup (\ordsucc{\beta}\times \{1\}), \rlexless}\\
	&= \ordtype{(\alpha\times \{0\})
	\cup (\beta \times \{1\}), \rlexless} \ordplus 1\\
	&= (\alpha \ordplus  \beta) \ordplus  1
\end{align*}
\footnote{Note de l'édition française : l'original remplace
à tort $0\times\{1\}=\emptyset$ par $\{0\}$ à la
deuxième ligne ; l'ensemble vide est rétabli ici.}
Soit maintenant $\beta \neq \emptyset$ un ordinal limite.
Si $\delta < \beta$, alors
$\delta\ordplus 1 < \beta$ aussi ; ainsi,
$\alpha \ordplus \delta$ est un segment initial propre
de $\alpha \ordplus \beta$. Par conséquent,
$\alpha \ordplus \beta$ est une borne supérieure stricte
de $X = \Setabs{\alpha \ordplus \delta}{\delta < \beta}$.
De plus, si
$\alpha \leq \gamma < \alpha \ordplus \beta$,
alors il est clair que
$\gamma = \alpha \ordplus \delta$ pour un
$\delta < \beta$. Donc
$\alpha \ordplus \beta =
\supstrict_{\delta< \beta}(\alpha\ordplus \delta)$.
\end{proof}

Voici cependant un fait frappant. Pour définir
l'addition ordinale, nous aurions \emph{aussi} pu
employer simplement le théorème de récursion transfinie
et poser les équations de récursion exactement comme
dans le \olref{ordadditionrecursion} (en écrivant
«~$\ordsucc{\beta}$~» plutôt que
«~$\beta \ordplus 1$~»).

Il existe donc deux manières de définir des opérations
sur les ordinaux. Nous pouvons les définir
\emph{synthétiquement}, en construisant explicitement
un ensemble bien ordonné et en prenant son type d'ordre.
Ou bien nous pouvons les définir \emph{récursivement},
en posant les équations de récursion. Si nous procédons
correctement, le résultat est identique :
Le \olref[ordinals][ordtype]{thmOrdinalRepresentation}
garantit que la construction synthétique donne un
ordinal, tandis que le théorème de récursion transfinie
garantit que les équations déterminent des ordinaux
\emph{uniques}.\footnote{Note de l'édition française :
l'original attribue l'unicité des solutions des équations
au théorème de représentation des ordinaux ; elle relève
du théorème de récursion transfinie.}

À bien des égards, l'arithmétique ordinale ressemble
à l'addition des nombres naturels. Nous pouvons,
par exemple, démontrer ceci :

\begin{lem}\ollabel{ordinaladditionisnice}
Si $\alpha, \beta, \gamma$ sont des ordinaux, alors :
\begin{enumerate}
	\item\ollabel{ordaddition1} si $\beta < \gamma$,
	alors $\alpha \ordplus \beta <
	\alpha \ordplus \gamma$ ;
	\item\ollabel{ordaddition2} si
	$\alpha \ordplus \beta =
	\alpha\ordplus \gamma$, alors
	$\beta = \gamma$ ;
	\item\ollabel{ordaddition3}
	$\alpha \ordplus (\beta \ordplus \gamma)
	= (\alpha \ordplus \beta) \ordplus \gamma$,
	c'est-à-dire que l'addition est associative ;
	\item\ollabel{ordaddition4} si
	$\alpha \leq \beta$, alors
	$\alpha \ordplus \gamma
	\leq \beta \ordplus \gamma$.
\end{enumerate}
\end{lem}

\begin{proof}
Démontrons \olref{ordaddition3} et laissons les autres
assertions en exercice. Raisonnons par récurrence
transfinie simple sur $\gamma$, en utilisant
le \olref{ordadditionrecursion}. Si $\gamma = 0$ :
\[
(\alpha \ordplus \beta) \ordplus 0
= \alpha \ordplus \beta
= \alpha \ordplus (\beta \ordplus 0)
\]
Si $\gamma = \delta\ordplus 1$, supposons, comme
hypothèse de récurrence, que
$(\alpha \ordplus \beta) \ordplus \delta =
\alpha \ordplus (\beta \ordplus \delta)$.
En appliquant trois fois le \olref{ordadditionrecursion} :
\begin{align*}
	(\alpha \ordplus \beta) \ordplus (\delta \ordplus 1)
	&= ((\alpha \ordplus \beta) \ordplus \delta)
	\ordplus 1\\
	&= (\alpha \ordplus (\beta \ordplus \delta))
	\ordplus 1\\
	&= \alpha \ordplus ((\beta \ordplus \delta)
	\ordplus 1)\\
	&= \alpha \ordplus (\beta \ordplus
	(\delta\ordplus 1))
\end{align*}
Si $\gamma$ est un ordinal limite, supposons que,
pour tout $\delta \in \gamma$,
$(\alpha \ordplus \beta) \ordplus \delta =
\alpha \ordplus (\beta \ordplus \delta)$.
Alors :
\begin{align*}
	(\alpha \ordplus \beta) \ordplus \gamma
	&= \supstrict_{\delta < \gamma}
	((\alpha \ordplus \beta) \ordplus \delta)\\
	&= \supstrict_{\delta < \gamma}
	(\alpha \ordplus (\beta \ordplus \delta))\\
	&= \alpha \ordplus \supstrict_{\delta < \gamma}
	(\beta \ordplus \delta)\\
	&= \alpha \ordplus (\beta \ordplus \gamma)
\end{align*}
\end{proof}

\begin{prob}
Démontrez les autres assertions du
\olref[sth][ord-arithmetic][add]{ordinaladditionisnice}.
\end{prob}

Sous ces aspects, l'addition ordinale nous est familière.
Mais, sur un point décisif, elle ne ressemble
\emph{pas} à l'addition des nombres naturels.

\begin{prop}\ollabel{ordsumnotcommute}
L'addition ordinale n'est {pas} commutative :
$1 \ordplus \omega = \omega <
\omega \ordplus 1$.
\end{prop}

\begin{proof}
Remarquons que
$1 \ordplus \omega =
\supstrict_{n < \omega} (1 \ordplus n)
= \omega \in \omega \cup \{\omega\}
= \ordsucc{\omega} = \omega \ordplus 1$.
\end{proof}
\noindent
Cela peut surprendre de prime abord, mais
\emph{cela ne devrait pas}. Pour
$1 \ordplus \omega$, nous considérons le type
d'ordre obtenu en plaçant un élément supplémentaire
\emph{avant} tous les nombres naturels. Comme dans
notre raisonnement sur l'hôtel de Hilbert dans la
\olref[sfr][infinite][hilbert]{sec}, ce premier
élément supplémentaire ne devrait intuitivement
pas changer le type d'ordre global. En revanche,
pour $\omega \ordplus 1$, l'élément supplémentaire
vient \emph{après} tous les nombres naturels.
Et c'est une situation radicalement différente !

\end{document}
